
\clearpage
\begin{proofbox}
\textbf{证明前的五步认知桥梁｜不是第四套孤立公式，而是同一条思路。}
\par\medskip\noindent
\textbf{第 1 步｜ 初中垂径定理：}
\par\nopagebreak\noindent
在圆中，半径相等且 $AM=MB$，两三角形全等，故非直径弦满足 $OM\perp AB$。
\par\medskip\noindent
\textbf{第 2 步｜ 圆的斜率：}
\par\nopagebreak\noindent
垂直关系给 $k=-x_0/y_0$（此处 $x_0y_0\ne0$）。
\par\medskip\noindent
\textbf{第 3 步｜ 椭圆拉伸：}
\par\nopagebreak\noindent
单位圆经 $x=au,\ y=bv$ 变为椭圆，
$k_{xy}=(b/a)k_{uv}$，从而 $k=-b^2x_0/(a^2y_0)$。
\par\medskip\noindent
\textbf{第 4 步｜ 双曲线变号：}
\par\nopagebreak\noindent
将方程中 $y^2/b^2$ 项系数由 $+1$ 改为 $-1$，
其对应中点弦斜率符号变为正；\textbf{这一步不是由实数拉伸推出的，必须证明。}
\par\medskip\noindent
\textbf{第 5 步｜ 点差统一：}
\par\nopagebreak\noindent
对
$\frac{x^2}{a^2}+\varepsilon\frac{y^2}{b^2}=1$
（$\varepsilon=1$ 为椭圆，$\varepsilon=-1$ 为双曲线），两端点相减，得
\[
\frac{x_0\Delta x}{a^2}+\varepsilon\frac{y_0\Delta y}{b^2}=0
\ \Longrightarrow\
\boxed{k=-\varepsilon\frac{b^2x_0}{a^2y_0}}.
\]
取 $a=b=R,\varepsilon=1$ 还原圆的斜率结论。
\par\medskip\noindent
\textbf{要点：}五步解释“为什么同源、哪里不同”；下列\textbf{方法一}最快、\textbf{方法二}门槛最低、\textbf{方法三}常规，分别解决真实的双曲线问题。
\end{proofbox}
\clearpage
\begin{proofbox}
\textbf{方法一｜最快证明：两点代入，直接相减（考场优先，三行）。}
\begin{center}\begin{tikzpicture}[x=1.1cm,y=1.1cm,font=\small,>=stealth]
\node[font=\small\bfseries] at (0,2.8) {点差法：横纵坐标差确定弦方向};
\draw[gray,->] (-4.6,0)--(4.7,0) node[right] {$x$};
\draw[gray,->] (0,-2.05)--(0,2.1) node[above] {$y$};
\node[below left] at (0,0) {$O$};
\draw[CStatement,thick,domain=-1.8:1.8,samples=100] plot ({2*sqrt(1+\x*\x)},\x);
\draw[CStatement,thick,domain=-1.8:1.8,samples=100] plot ({-2*sqrt(1+\x*\x)},\x);
\coordinate (A) at (10/3,4/3); \coordinate (B) at (-2,0); \coordinate (M) at (2/3,2/3);
\draw[teal,very thick] (B)--(A);
\fill (A) circle(1.8pt) node[above left] {$A$};
\fill (B) circle(1.8pt) node[below left] {$B$};
\fill[orange!80!black] (M) circle(2pt) node[below right] {$M$};
\draw[purple,dashed,thick] (B)--(10/3,0)--(A);
\node[below] at (2/3,-.5) {$x_1-x_2$};
\node[right] at (3.65,.65) {$y_1-y_2$};

\node[align=center] at (0,-2.65) {$M=\bigl((x_1+x_2)/2,(y_1+y_2)/2\bigr)$；图取 $a=2,b=1$};
\end{tikzpicture}
\end{center}
设 $A(x_1,y_1)$、$B(x_2,y_2)$ 是两个不同交点，$M(x_0,y_0)$ 为其中点。由
\[
\frac{x_1^2}{a^2}-\frac{y_1^2}{b^2}=1,\qquad
\frac{x_2^2}{a^2}-\frac{y_2^2}{b^2}=1
\]
相减，并用 $x_1+x_2=2x_0,\ y_1+y_2=2y_0$，得
\[
\underbrace{\frac{2x_0(x_1-x_2)}{a^2}}_{\text{横坐标平方差}}
=\underbrace{\frac{2y_0(y_1-y_2)}{b^2}}_{\text{纵坐标平方差}}.
\]
若 $y_0\ne0$ 且 $x_1=x_2$，则上式推出 $y_1=y_2$，矛盾；因此
\[
\boxed{k=\frac{y_1-y_2}{x_1-x_2}=\frac{b^2x_0}{a^2y_0}}.
\]
\textbf{为什么最快：}不设直线、不联立求交点，直接利用``平方差 = 差$\times$和''。
\textbf{定位：}此法只回答``往哪里走''；题目若问\textbf{能否构成实际中点弦}，
还要用后面的\textbf{相加位移关系}或判别式判定。
\end{proofbox}
\clearpage
\begin{proofbox}
\textbf{方法二｜最低门槛：把两个端点关于中点对称，再看``方向 + 距离''。}
\par\smallskip\noindent
\textbf{先接上初中圆的经验：}圆的垂径定理给方向 $OM\perp PQ$，勾股定理给半弦长度
$OM^2+MP^2=R^2$。这是两件不同的事：\textbf{方向正确，还要走得够远，才会到达曲线。}
双曲线遵循相同的解题逻辑，但``勾股关系''变成\textbf{带符号的二次型等式}。
\begin{center}\begin{tikzpicture}[x=1.1cm,y=1.1cm,font=\small,>=stealth]
\node[font=\small\bfseries] at (0,2.8) {对称法：两端点从中点作相反位移};
\draw[gray,->] (-4.6,0)--(4.7,0) node[right] {$x$};
\draw[gray,->] (0,-2.05)--(0,2.1) node[above] {$y$};
\node[below left] at (0,0) {$O$};
\draw[CStatement,thick,domain=-1.8:1.8,samples=100] plot ({2*sqrt(1+\x*\x)},\x);
\draw[CStatement,thick,domain=-1.8:1.8,samples=100] plot ({-2*sqrt(1+\x*\x)},\x);
\coordinate (A) at (10/3,4/3); \coordinate (B) at (-2,0); \coordinate (M) at (2/3,2/3);
\draw[teal,very thick] (B)--(A);
\fill (A) circle(1.8pt) node[above left] {$A$};
\fill (B) circle(1.8pt) node[below left] {$B$};
\fill[orange!80!black] (M) circle(2pt) node[below right] {$M$};
\draw[purple,->,thick,shorten <=3pt,shorten >=3pt] (M)--(A);
\draw[orange!80!black,->,thick,shorten <=3pt,shorten >=3pt] (M)--(B);
\node at (1.5,1.7) {$(u,v)$};
\node at (-1,-.7) {$(-u,-v)$};

\node[align=center] at (0,-2.65) {$A=M+(u,v),\quad B=M-(u,v)$；图取 $a=2,b=1$};
\end{tikzpicture}
\end{center}
\textbf{第 1 步｜中点条件先写进坐标。}设两个端点为
\[
\boxed{P=(x_0+u,y_0+v),\quad Q=(x_0-u,y_0-v)}.
\]
任取 $u,v$，其中点都是 $M(x_0,y_0)$；但端点能否\textbf{真在双曲线上}尚未确定。
\textbf{第 2 步｜让两个端点同时满足原方程（完整保留 $a,b$）。}
\[
\begin{aligned}
\frac{(x_0+u)^2}{a^2}-\frac{(y_0+v)^2}{b^2}&=1\qquad\text{(1)}\\[3pt]
\frac{(x_0-u)^2}{a^2}-\frac{(y_0-v)^2}{b^2}&=1\qquad\text{(2)}
\end{aligned}
\]
\textbf{第 3 步｜相减：去掉平方的公共部分，确定半弦方向。}
\[
(1)-(2):\quad \frac{4x_0u}{a^2}-\frac{4y_0v}{b^2}=0
\quad\Longrightarrow\quad
\boxed{\frac{x_0u}{a^2}=\frac{y_0v}{b^2}}.\qquad\text{(3)}
\]
此处只说明\textbf{若有真实端点，$(u,v)$ 必须沿这个方向}；并未保证能走到曲线上。
当 $y_0\ne0$ 且端点不同，得到 $k_{PQ}=v/u=b^2x_0/(a^2y_0)$。
\par\smallskip\noindent
\textbf{第 4 步｜为什么自然想到相加？}
两个端点的坐标是``$+u,+v$''与``$-u,-v$''：
\[
\begin{aligned}
(x_0+u)^2+(x_0-u)^2&=2x_0^2+2u^2,\\
(y_0+v)^2+(y_0-v)^2&=2y_0^2+2v^2.
\end{aligned}
\]
\textbf{相加恰好消去交叉项}，把中点的位置与偏移大小分离：
\[
\begin{aligned}
(1)+(2):\quad
\frac{2x_0^2+2u^2}{a^2}-\frac{2y_0^2+2v^2}{b^2}&=2,\\
\boxed{\underbrace{\left(\frac{x_0^2}{a^2}-\frac{y_0^2}{b^2}\right)}_{S}
+\underbrace{\left(\frac{u^2}{a^2}-\frac{v^2}{b^2}\right)}_{\text{半弦位移的带符号二次型}}}&=1.\qquad\text{(4)}
\end{aligned}
\]
与圆的 $OM^2+MP^2=R^2$ 是\textbf{同一个``位置 + 位移''结构}，
但\textbf{双曲线第二项不是 $MP^2$，不能当普通勾股定理使用}。
\par\smallskip\noindent
\textbf{第 5 步｜把方向与位移条件合起来，才判断存在性。}
当 $M\ne O$ 时，式 (3) 的所有方向向量都可写成
\[
\boxed{u=a^2y_0t,\quad v=b^2x_0t}\qquad(t\in\mathbb R).
\]
将其代入式 (4)，不消去 $a^2,b^2$：
\[
\begin{aligned}
S+\frac{(a^2y_0t)^2}{a^2}-\frac{(b^2x_0t)^2}{b^2}&=1,\\
S+(a^2y_0^2-b^2x_0^2)t^2&=1,\\
S-a^2b^2St^2&=1,\\
\boxed{t^2=\frac{S-1}{a^2b^2S}}&\quad(S\ne0).
\end{aligned}
\]
由于 $a^2b^2>0$，非退化实弦需要 $t\ne0$ 且为实数，即
\[
\boxed{t^2>0\iff \frac{S-1}{S}>0\iff S<0\ \text{或}\ S>1}.
\]
具体地，$0<S<1$ 时 $t^2<0$，方向虽然算得出却\textbf{没有真实位移}；
$S=1$ 时 $t=0$、两端点重合；$S=0,M\ne O$ 时 $0=1$，无解。
$M=O$ 不适用唯一候选方向的设法，单独取 $P=(a,0)$、$Q=(-a,0)$ 即得中心弦。
\par\smallskip\noindent
\textbf{反向充分性：}若 $S<0$ 或 $S>1$，可取 $t=\sqrt{(S-1)/(a^2b^2S)}>0$，
定义上式的 $u,v$ 和 $P,Q$；式 (3)、(4) 同时成立，反推出 (1)、(2)，
因此端点\textbf{确实存在且不同}。
\textbf{为什么门槛最低：}中点对称 $\to$ 相减\textbf{方向} $\to$ 相加\textbf{位移} $\to$ 实数平方判定；
这和初中``垂径定理定方向、勾股定理定长度''是一条认知主线。
\end{proofbox}
\clearpage
\begin{proofbox}
\textbf{方法三｜常规证明：设直线，联立二次方程，使用韦达。}
\begin{center}\begin{tikzpicture}[x=1.1cm,y=1.1cm,font=\small,>=stealth]
\node[font=\small\bfseries] at (0,2.8) {韦达法：中点投影使两根互为相反数};
\draw[gray,->] (-4.6,0)--(4.7,0) node[right] {$x$};
\draw[gray,->] (0,-2.05)--(0,2.1) node[above] {$y$};
\node[below left] at (0,0) {$O$};
\draw[CStatement,thick,domain=-1.8:1.8,samples=100] plot ({2*sqrt(1+\x*\x)},\x);
\draw[CStatement,thick,domain=-1.8:1.8,samples=100] plot ({-2*sqrt(1+\x*\x)},\x);
\coordinate (A) at (10/3,4/3); \coordinate (B) at (-2,0); \coordinate (M) at (2/3,2/3);
\draw[teal,very thick] (B)--(A);
\fill (A) circle(1.8pt) node[above left] {$A$};
\fill (B) circle(1.8pt) node[below left] {$B$};
\fill[orange!80!black] (M) circle(2pt) node[below right] {$M$};
\draw[purple,dashed] (10/3,4/3)--(10/3,-.8);
\draw[purple,dashed] (2/3,2/3)--(2/3,-.8);
\draw[purple,dashed] (-2,0)--(-2,-.8);
\draw[<->,purple] (-2,-1)--(2/3,-1);
\draw[<->,purple] (2/3,-1)--(10/3,-1);
\node[below] at (-2/3,-1.1) {$-h$};
\node[below] at (2,-1.1) {$h$};

\node[align=center] at (0,-2.65) {$t=x-x_0:\quad t_B=-h,\quad t_A=h,\quad t_A+t_B=0$};
\end{tikzpicture}
\end{center}
（讨论 $y_0\ne0$，已知弦确有两个不同交点。）设
$y=y_0+k(x-x_0)$，令 $t=x-x_0$，代入得
\[
\left(\frac1{a^2}-\frac{k^2}{b^2}\right)t^2
+2\left(\frac{x_0}{a^2}-\frac{ky_0}{b^2}\right)t
+(S-1)=0.
\]
弦上两个交点的 $t_1,t_2$ 满足 $t_1+t_2=0$。两交点不同且 $y_0\ne0$，所以对应非退化二次方程（若二次项为零，则 $k=\pm b/a$；再令一次项为零，会推出 $S=0$，常数项却为 $-1$，故不可能恒等为零）。韦达给出
\[
\frac{x_0}{a^2}-\frac{ky_0}{b^2}=0
\ \Longrightarrow\ \boxed{k=\frac{b^2x_0}{a^2y_0}}.
\]
反过来令 $k$ 取此候选值；$M\ne O$、$S\ne0$ 时，关于 $t$ 的方程化为
\[
-\frac{b^2S}{a^2y_0^2}t^2+S-1=0,
\quad\text{即}\quad
t^2=\frac{a^2y_0^2(S-1)}{b^2S}.
\]
实交点且不重合仍当且仅当 $S<0$ 或 $S>1$。$S=0,M\ne O$ 时原方程成为 $-1=0$；$M=O$ 则须另按中心弦方向讨论，见末节。
\textbf{为什么保留：}与教材的联立—韦达—判别式体系完全接轨，适合题目已给直线方程的情形，但运算更长。
\end{proofbox}
\clearpage
\begin{proofbox}
\textbf{补充方法四｜统一方程法：直接从中点与方向的共轭关系读出弦。}
\begin{center}\begin{tikzpicture}[x=1.1cm,y=1.1cm,font=\small,>=stealth]
\node[font=\small\bfseries] at (0,2.8) {统一方程法：反射的是两个曲线端点};
\draw[gray,->] (-4.6,0)--(4.7,0) node[right] {$x$};
\draw[gray,->] (0,-2.05)--(0,2.1) node[above] {$y$};
\node[below left] at (0,0) {$O$};
\draw[CStatement,thick,domain=-1.8:1.8,samples=100] plot ({2*sqrt(1+\x*\x)},\x);
\draw[CStatement,thick,domain=-1.8:1.8,samples=100] plot ({-2*sqrt(1+\x*\x)},\x);
\coordinate (A) at (10/3,4/3); \coordinate (B) at (-2,0); \coordinate (M) at (2/3,2/3);
\draw[teal,very thick] (B)--(A);
\fill (A) circle(1.8pt) node[above left] {$P$};
\fill (B) circle(1.8pt) node[below right,yshift=-8pt] {$2M-P$};
\fill[orange!80!black] (M) circle(2pt) node[below right] {$M$};
\draw[purple,dashed,thick] (-2.4,-.1)--(3.7,1.425);

\node[align=center] at (0,-2.65) {$P+(2M-P)=2M$：反射端点确定弦线；图取 $a=2,b=1$};
\end{tikzpicture}
\end{center}
设 $P=(x,y)$ 是弦的一个端点，$M\ne O$ 为中点，则 $2M-P$ 是弦的另一端点。
两点的双曲线方程相减：
\[
\frac{x^2-(2x_0-x)^2}{a^2}-\frac{y^2-(2y_0-y)^2}{b^2}=0
\quad\Longleftrightarrow\quad
\boxed{\frac{x_0x}{a^2}-\frac{y_0y}{b^2}=S}.
\]
由于两个端点都满足该一次方程，且 $M\ne O$ 使其系数不全为零，它确定的直线就是整条弦所在直线。这里不能将弦上任意点直接当作曲线端点。
\textbf{为什么值得补充：}一次涵盖斜弦、水平弦、竖直弦；无需先计算斜率，适合求弦方程及轨迹题。
\end{proofbox}
\clearpage
\begin{proofbox}
\textbf{严格补全｜异支、同支及中心例外。}
\begin{center}\begin{tikzpicture}[x=.62cm,y=.62cm,font=\small,>=stealth]
\node[font=\small\bfseries] at (0,4.6) {同一支：无交点、相切、竖直弦};
\draw[CStatement,thick,domain=-1.8:1.8,samples=100] plot ({2*sqrt(1+\x*\x)},\x);
\draw[CStatement,thick,domain=-1.8:1.8,samples=100] plot ({-2*sqrt(1+\x*\x)},\x);
\draw[gray,->] (-4.5,0)--(4.6,0);
\draw[CTrap,dotted,thick] (1,-1.8)--(1,1.8);
\draw[purple,dashed,thick] (2,-1.8)--(2,1.8);
\draw[teal,very thick] (3,-1.11803399)--(3,1.11803399);
\fill (2,0) circle(2pt);
\fill[orange!80!black] (3,0) circle(2pt);
\node[above] at (.9,2.15) {$x=1$};
\node[above] at (2,2.8) {$x=2$};
\node[above] at (3.6,2.15) {$x=3$};
\node[align=center] at (0,-2.9) {$S=1$ 仅有切点\\$S>1$ 才有竖直弦};
\begin{scope}[xshift=6.8cm]
\node[font=\small\bfseries] at (0,4.6) {中心例外：只有部分方向可取};
\draw[CStatement,thick,domain=-1.8:1.8,samples=100] plot ({2*sqrt(1+\x*\x)},\x);
\draw[CStatement,thick,domain=-1.8:1.8,samples=100] plot ({-2*sqrt(1+\x*\x)},\x);
\draw[gray,->] (-4.5,0)--(4.6,0);
\draw[purple,dashed] (-3.8,-1.9)--(3.8,1.9);
\draw[teal,very thick] (-2.5,-.75)--(2.5,.75);
\fill[orange!80!black] (0,0) circle(2pt); \node[orange!80!black] at (0,-1.35) {$M=O$};
\node[align=center] at (0,2.7) {虚线是渐近线：$k=1/2$};
\node[align=center] at (0,-2.9) {实线中心弦：$k=3/10$\\存在范围 $|k|<1/2$};
\end{scope}
\node at (5.5,-4.25) {两图均取 $x^2/4-y^2=1$，相同坐标尺度};
\end{tikzpicture}
\end{center}
对 $M\ne O$ 且弦存在的端点 $A=M+t(a^2y_0,b^2x_0)$、$B=M-t(a^2y_0,b^2x_0)$，有
\[
x_Ax_B=x_0^2-a^4y_0^2t^2
=a^2S+\frac{a^2y_0^2}{b^2S}
=\frac{a^2(S^2+y_0^2/b^2)}{S}.
\]
因 $a^2(S^2+y_0^2/b^2)>0$，所以 $x_Ax_B$ 与 $S$ 同号：$S<0$ 为异支，$S>1$ 为同支。
若 $M=O$，两点取 $(\pm x,\pm y)$，任意 $y=kx$ 代入得
\[
x^2\left(\frac1{a^2}-\frac{k^2}{b^2}\right)=1,
\]
因而中心弦存在当且仅当 $|k|<b/a$（斜率 $\pm b/a$ 恰是渐近线，不能取）。
至此，对所有中点 $M$ 的完整存在性判定为 $S<0$ 或 $S>1$ 或 $M=O$。
\end{proofbox}
